% ============================================================
%  B 组实验操作单（正文版，供 \input 进 menu.tex / 独立打印共用）
%  语言刻意简化，供任何人照着执行；日期以 2026-09-23 为第 1 天
%  若延期开工，把文中日期整体后移即可
% ============================================================

\section{B 组实验操作单（2026-09-23 起，照着做）}
\label{sec:Bop}

{\small 本操作单供他人独立执行；日期按 09-23 起算，若延期开工，把下面的日期整体后移即可。每一步都写了"做对了的样子"，对不上就停下报告。}

\subsection*{一、开工前，先把这些放在桌上（放好了打勾）}
\begin{itemize}
  \item （\hspace{1.2em}）台面上有三个烧杯，贴着标签：\textbf{基液}、\textbf{外层液}、\textbf{中层液}（以前配好、还没用完的）
  \item （\hspace{1.2em}）正丙醇锆（小瓶，怕水，用完马上盖紧）
  \item （\hspace{1.2em}）葡萄糖（白色粉末）2.00\,g；去离子水 3.0\,mL
  \item （\hspace{1.2em}）天平（能称到 0.01\,g）、玻璃棒、滴管或注射器
  \item （\hspace{1.2em}）三个 10\,mL 注射器 + \textbf{全金属针头}（塑料针头会被液体泡坏）
  \item （\hspace{1.2em}）硅油纸（约 30\,cm 长条）、胶带
  \item （\hspace{1.2em}）静电纺丝机、马弗炉、管式炉（通氩气那台）
  \item （\hspace{1.2em}）笔和本子（第五节要填数字）
\end{itemize}

\subsection*{二、时间表：什么时候干什么}

\begin{table}[htbp]
\centering
\caption{B 组操作时间表（09-23 起）}
\begin{tabularx}{\textwidth}{p{2.7cm} X}
\toprule
\textbf{时间} & \textbf{要做的事} \\
\midrule
\textbf{09-23（三）}\newline 09:00 & \textbf{（1）先检查液体。}把三个烧杯举起来看一看、用玻璃棒挑一挑。\textbf{对的样子}：透明、不浑、没有白色絮状、没有胶块，挑起来是短短一小段丝。\textbf{如果发浑、有胶块、或一挑就拉长丝不断，立刻停下}，告诉负责人（这批液体放了 18 天，可能坏了）。\\
\midrule
\textbf{09-23（三）}\newline 09:20 & \textbf{（2）试纺。}倒 2--3\,mL 出来，在小纸片上纺一下，看喷出来的丝稳不稳（不停不断、不滴液）。稳了再往下做。\\
\midrule
\textbf{09-23（三）}\newline 09:40 & \textbf{（3）加锆。}用注射器往\textbf{中层液}里慢慢滴入正丙醇锆 \textbf{0.64\,g}，一边滴一边搅，滴完继续搅 \textbf{30 分钟}。\textbf{对的样子}：液体还是透明，没有变浑、没有白点；\textbf{有白点就停下报告}。\\
\midrule
\textbf{09-23（三）}\newline 10:20 & \textbf{（4）配糖水。}葡萄糖 2.00\,g + 去离子水 3.0\,mL，隔温水（40--50\,°C）搅到完全化开、没有颗粒，再放到不烫手。\\
\midrule
\textbf{09-23（三）}\newline 10:40 & \textbf{（5）把糖水倒进中层液}，搅 \textbf{15--30 分钟}。\textbf{【最重要】搅完 30 分钟以内必须开始纺丝}，不能放着过夜。\\
\midrule
\textbf{09-23（三）}\newline 11:10 起\newline 到 09-24（四）早上 & \textbf{（6）纺丝，约 20 小时，中间不能停。}顺序：先纺\textbf{外层液} 10\,mL，再纺\textbf{中层液} 10\,mL，最后再纺\textbf{外层液} 10\,mL。出液速度 1.5\,mL/h（机器设定好不要改），硅油纸接住。每层开始前先空喷一小会儿，等丝稳定了再让硅油纸接。晚上没人也要开着；中途看一眼，正常就是细细的丝不断。\\
\midrule
\textbf{09-24（四）}\newline 07:00 & \textbf{（7）取膜。}把膜连着硅油纸拿下来，平放，称一次重量（\textbf{初重}）记下来，晾 1--2 小时。\textbf{不要折叠。}\\
\midrule
\textbf{09-24（四）}\newline 14:00 & \textbf{（8）空气预氧化（马弗炉，不要盖东西）。}室温慢慢升到 \textbf{100\,°C} 停 30 分钟；再升到 \textbf{150\,°C} 停 30 分钟；再升到 \textbf{220\,°C} 停 60 分钟；然后关火，等炉子自己凉到 \textbf{60\,°C 以下}再开门取出。\textbf{全过程不能超过 220\,°C。}\textbf{对的样子}：膜变成均匀的\textbf{金黄色或浅褐色}、还是完整一片；称一次重量记下（预氧化后重）。\\
\midrule
\textbf{09-24（四）}\newline 18:00 & \textbf{（9）进管式炉（通氩气）。}把膜平放进瓷舟（写上编号），放进管式炉，\textbf{先通氩气 30 分钟}把空气赶走，然后开程序：\textbf{200 分钟}升到 \textbf{200\,°C} 停 60 分钟；\textbf{400 分钟}升到 \textbf{900\,°C} 停 60 分钟；关火，让它自己慢慢凉（整夜）。\\
\midrule
\textbf{09-25（五）}\newline 早上 & \textbf{（10）出炉。}等炉子凉到 \textbf{100\,°C 以下}再取。称一次重量（\textbf{终重}）记下来；装进样品袋，写标签：\textbf{B + 日期 + 名字}。\\
\midrule
\textbf{09-25（五）}\newline 上午 & \textbf{（11）拍照 + 送检。}拍两张照片（整片的样子、用手弯一下的样子）；联系送 \textbf{SEM}，说明"要看表面，能切就再看截面"。\\
\bottomrule
\end{tabularx}
\end{table}

\subsection*{三、绝对不能做错的 6 件事}
\begin{enumerate}
  \item \textbf{顺序不能反}：先加锆，后加糖水；反过来锆会析出白点。
  \item \textbf{糖水加进去 30 分钟内必须开始纺丝}，不能放。
  \item \textbf{针头必须全金属}，塑料的会被液体泡坏、堵住。
  \item \textbf{预氧化不能超过 220\,°C}，升温也不能快，否则膜会烧黑、发脆。
  \item \textbf{两个炉子的顺序不能反}：先马弗炉（有空气），后管式炉（通氩气）；换炉子前必须先凉到 60\,°C 以下。
  \item \textbf{管式炉一定要先通氩气 30 分钟}，把里面的空气赶走再升温。
\end{enumerate}

\subsection*{四、出问题了怎么办}
\begin{itemize}
  \item 加锆后\textbf{出现白点、变浑} $\rightarrow$ 马上停下，不要纺，记下来报告；
  \item 纺丝时\textbf{老是断、老是滴} $\rightarrow$ 检查针头堵没堵，换个针头再试；
  \item 纺出来的膜\textbf{太薄、像没纺上} $\rightarrow$ 叫负责人来看，先别关机器；
  \item 预氧化后\textbf{颜色发黑、一碰就碎} $\rightarrow$ 说明温度高了，记下来，下次降到 200\,°C；
  \item 液体\textbf{变得很稠、一挑就是长丝} $\rightarrow$ 记下时间，不要再硬纺，报告负责人（可能要重配液体）。
\end{itemize}

\subsection*{五、要填的数字（做完一项填一项）}
\begin{table}[htbp]
\centering
\caption{操作记录（操作者填写）}
\begin{tabularx}{\textwidth}{X p{3.2cm} p{3.2cm}}
\toprule
\textbf{项目} & \textbf{要求} & \textbf{实际} \\
\midrule
三个烧杯的检查结果 & 透明 / 不透明 & \\
试纺是否稳定 & 稳定 / 不稳 & \\
正丙醇锆称了多少克 & 0.64\,g & \\
葡萄糖 / 水的量 & 2.00\,g / 3.0\,mL & \\
搅完糖水到开始纺丝隔多久 & 30 分钟以内 & \\
纺丝开始时间 / 结束时间 & —— & \\
初重（纺完） & —— & \\
预氧化后的颜色 & 金色 / 褐色 & \\
预氧化后重量 & —— & \\
终重（烧完） & —— & \\
\bottomrule
\end{tabularx}
\end{table}

\noindent\textbf{做完把这张操作单和第五节填好的数字一起交给负责人。}
